\documentclass[11pt,a4paper]{article}

% --- PAQUETS DE BASE ---
\usepackage[utf8]{inputenc}
\usepackage[T1]{fontenc}
\usepackage[french]{babel}
\usepackage{amsmath, amsfonts, amssymb}
\usepackage{geometry}
\geometry{top=2cm, bottom=2.5cm, left=2cm, right=2cm}

% --- DESIGN ---
\usepackage{xcolor}
\usepackage{colortbl}
\definecolor{BrandPrimary}{HTML}{0D9488}
\definecolor{BrandDark}{HTML}{111827}

\usepackage{tcolorbox}
\tcbuselibrary{skins, breakable}

\usepackage{titlesec}
\titleformat{\section}{\color{BrandPrimary}\normalfont\Large\bfseries}{\thesection}{1em}{}[{\titlerule[0.8pt]}]

% --- POLICE ---
\usepackage{helvet}
\renewcommand{\familydefault}{\sfdefault}

\begin{document}

\begin{tcolorbox}[enhanced, colback=BrandDark, colframe=BrandPrimary, arc=10pt, boxrule=2pt, halign=center]
    \vspace{0.3cm}
    {\Huge \bfseries \color{white} ZENITH LEARN} \\
    \vspace{0.2cm}
    {\Large \color{BrandPrimary} IA Éthique : Modélisation TWIST 03} \\
    \vspace{0.3cm}
\end{tcolorbox}

\section{Analyse sémantique : ``Court vs Utile''}
Le \textbf{TWIST 03} dénonce le biais de complétion : \textit{``Les cours les plus complétés sont les plus courts, pas les plus utiles.''} L'IA doit neutraliser ce biais mécanique pour ne pas enterrer les cours de fond.

Ce phénomène est bien documenté dans les sciences de l'éducation : face à deux formations, l'une de 30 minutes et l'autre de 8 heures, l'algorithme naïf verra toujours la première avec un meilleur taux de complétion, même si la seconde est objectivement plus structurante. Le TWIST 03 corrige ce biais en donnant au taux de complétion un \textbf{poids proportionnel à l'effort consenti}.

\section{Modélisation Mathématique du Redressement}
Nous cherchons une fonction de mérite qui valorise l'effort consenti en fonction de la durée du contenu.

\subsection{Formule de la Complétion Ajustée ($S_C^{adj}$)}
\begin{equation}
S_C^{adj} = C_{raw} \cdot \Theta(D)
\end{equation}

\subsection{Genèse de la formule : pourquoi cette construction ?}

Le point de départ est un constat d'injustice : le taux brut $C_{raw}$ ne dit rien sur l'effort fourni. Finir un cours de 20 minutes à 100\% est un acte trivial ; finir un cours de 40 heures à 100\% est un exploit. Pour corriger cela, on multiplie le taux brut par un \textbf{facteur de pondération $\Theta(D)$} qui augmente avec la durée. L'enjeu est de choisir une fonction mathématique qui croît avec $D$ de manière réaliste : ni trop lentement (la correction serait insuffisante), ni trop vite (un cours de 100 heures ne devrait pas peser infiniment plus qu'un cours de 50 heures).

\subsection{Détail de la Fonction de Pondération $\Theta(D)$}
\begin{equation}
\Theta(D) = \frac{\ln(1 + D)}{\ln(1 + D_{med})}
\end{equation}

\begin{itemize}
    \item \textbf{$C_{raw}$ (Taux brut) :} La valeur de complétion brute, comprise entre 0 et 1. C'est la donnée d'entrée non corrigée, telle qu'elle serait utilisée dans un système naïf.

    \item \textbf{$D$ (Durée théorique) :} Exprimée en heures, c'est la durée totale du cours telle que définie par son créateur. C'est la variable qui permet à $\Theta$ de s'adapter à chaque cours individuellement.

    \item \textbf{Pourquoi le Logarithme ($\ln$) ?}
    Le logarithme est la fonction naturelle pour modéliser la \textbf{difficulté cognitive à rendement décroissant}. Voici l'intuition : passer d'un cours de 1 heure à un cours de 5 heures représente un saut d'engagement considérable (multiplier son investissement par 5 est un effort psychologique énorme). Passer d'un cours de 30 heures à 34 heures représente un ajout marginal (4 heures supplémentaires sur une base déjà massive changent peu la nature de l'effort). Le logarithme capture exactement cette asymétrie : il croît vite pour les petites valeurs de $D$ et ralentit pour les grandes, reflétant ainsi la réalité psychologique de l'engagement. Une fonction linéaire serait injuste dans l'autre sens : elle sur-récompenserait les cours très longs de manière disproportionnée.

    \item \textbf{Le terme $+1$ dans $\ln(1 + D)$ :} Ce décalage de 1 est une précaution technique. Mathématiquement, $\ln(0)$ est indéfini (tend vers $-\infty$). En ajoutant 1, on s'assure que même pour un cours de durée $D = 0$ (un contenu théoriquement instantané), la fonction reste définie et vaut $\ln(1) = 0$. C'est une ``garde-fou'' numérique qui rend la formule robuste aux cas limites.

    \item \textbf{$D_{med}$ (Durée Médiane) : Le point pivot.} C'est la valeur autour de laquelle l'ajustement est calibré. Si un cours a exactement la durée médiane du catalogue, alors $\ln(1 + D) = \ln(1 + D_{med})$, et donc $\Theta(D_{med}) = 1$~: l'IA ne modifie rien pour ce cours, il est à la norme. Si le cours est \textbf{plus long} que la médiane, $\Theta(D) > 1$~: il reçoit un ``Bonus de Mérite'', son taux de complétion est amplifié. Si le cours est \textbf{plus court} que la médiane, $\Theta(D) < 1$~: il reçoit un ``Malus de Facilité'', son taux de complétion est dégonflé. La durée médiane (et non la moyenne) est choisie comme pivot car elle est plus robuste aux valeurs aberrantes : un seul cours de 200 heures ne déplace pas la médiane de manière catastrophique, contrairement à la moyenne.
\end{itemize}

\section{Nouvelle Structure du Dataset (Spécification T03)}
\begin{table}[h!]
\centering
\begin{tabular}{|l|l|p{7cm}|}
\hline
\rowcolor{BrandDark} \color{white} Colonne & \rowcolor{BrandDark} \color{white} Type & \rowcolor{BrandDark} \color{white} Rôle Algorithmique \\ \hline
\texttt{duree\_prevue} & Float & Base logarithmique du redressement. Valeur de $D$ dans $\Theta$. \\ \hline
\rowcolor{BrandPrimary!20} \texttt{adj\_completion} & Float & Signal de succès académique ``Nettoyé''. Résultat de $S_C^{adj}$. \\ \hline
\end{tabular}
\end{table}

\section{Impact sur les acquis précédents}
\begin{itemize}
    \item \textbf{Évolution du TWIST 01} : La complétion académique n'est plus une valeur brute mais dynamique. Le 100\% d'un cours court pèse désormais moins que le 100\% d'un cours long. Le signal $S_C^{adj}$ vient s'ajouter au registre de signaux disponibles, aux côtés du $S_E$ du TWIST 01.
    \item \textbf{Consolidation du TWIST 02} : En neutralisant le biais de facilité, on décourage le comportement de ``chasse aux badges'' faciles. Un apprenant qui enchaîne des micro-cours pour gonfler son taux de complétion ne verra plus ses métriques s'envoler artificiellement.
\end{itemize}

\section{Conclusion}
Le TWIST 03 transforme la complétion d'une simple mesure de parcours en une mesure de \textbf{Poids Pédagogique}, indispensable pour valoriser les formations entrepreneuriales structurantes. Il introduit un principe fondamental qui traversera tous les twists suivants : \textbf{la valeur d'un signal doit toujours être contextualisée par l'effort qu'il représente}.

\end{document}